\documentclass[11pt,a4paper]{article}
\usepackage[margin=1in]{geometry}
\usepackage{hyperref}
\usepackage{enumitem}
\usepackage{graphicx}
\usepackage{xcolor}

% -----------------------------------------------------
% bvraghav-tiet-ucs503p-article.sty
% -----------------------------------------------------

% Variable: Lab Instructor
\makeatletter \newcommand{\BvrSetLabInstructor}[1]{%
  \gdef\Bvr@LabInstructor{#1}}
% default
\newcommand{\Bvr@LabInstructor}{Dr. Raghav %
  B. Venkataramaiyer}
% public accessor
\newcommand{\BvrLabInstructorValue}{\Bvr@LabInstructor}
\makeatother

% Add subtitle
\usepackage{titling}
% -- Set title bold
\pretitle{\begin{center}\bfseries\fontsize{18bp}{18bp}\selectfont}
\makeatletter
\newcommand{\BvrSubtitle}[1]{%
  \posttitle{%
    \par\end{center}
    \begin{center}\large#1\end{center}
    \vskip 0.5em}%
}
\makeatother

% Hints in the section title(s)
\newcommand{\BvrHint}[1]{%
  {\normalfont\normalsize\itshape\hfill #1}}

% -----------------------------------------------------

\title{Smart Personal Finance \& Goal Management System}

\BvrSetLabInstructor{Ms.Nisha Thakur}

\BvrSubtitle{%
  Project Proposal for UCS503P  \\
  Submitted to: \BvrLabInstructorValue \\ \bfseries
  Thapar Institute of Engineering and Technology%
}

\author{
  Riya Garg, CSED%
  \thanks{Roll No: \texttt{1024030094}}
  \and
  Niharika, CSED%
  \thanks{Roll No: \texttt{1024030450}}
  \and
  Anushka Aggarwal, CSED%
  \thanks{{Roll No: \texttt{1024030443}}
  \and
  Snehal Agarwal, CSED%
  \thanks{{Roll No: \texttt{1024030259}}
}

\date{\today}

\begin{document}
\maketitle

\section{Higher-order goal%
  \BvrHint{\ldots secondary framing}}

This project aims to improve personal financial decision-making by helping users
not only track their income and expenses, but also understand their financial
behavior and take concrete actions toward long-term goals. Most personal
finance applications primarily present financial information through balances,
transactions, and charts. However, users often struggle to translate this
information into consistent financial behavior.

The higher-order goal of this project is therefore to build a
\textbf{goal-oriented personal finance system} that converts financial data into
actionable recommendations, supports planning under uncertainty, and helps
users make informed decisions while maintaining privacy and transparency.


\section{Time-to-value%
  \BvrHint{\ldots primary heuristic}}

\begin{itemize}[leftmargin=0.7cm]
    \item We will deliver meaningful functionality incrementally, beginning with
    secure transaction tracking, financial goals, budgeting, and a basic
    dashboard.
    
    \item We will prioritize rapid feedback by developing the system in
    iterations and validating each major feature through automated tests and
    user-oriented scenarios.
    
    \item Advanced features such as behavioral pattern detection, anomaly
    detection, scenario simulation, and shared-finance workflows will be added
    incrementally after the core financial model is stable.
    
    \item CI/CD automation will be introduced early so that every significant
    change can be tested consistently and deployed in a repeatable manner.
\end{itemize}


\section{Problem Statement}

Many personal finance applications focus primarily on displaying financial
information such as balances, transactions, spending categories, and monthly
summaries. While these features are useful for understanding past spending,
they do not necessarily help users change their financial behavior or determine
whether they are on track to achieve future goals.

Common challenges include:

\begin{itemize}[leftmargin=0.7cm]
    \item \textbf{Limited goal orientation:} users can see their spending but
    may not understand how their current behavior affects goals such as
    building an emergency fund, purchasing a laptop, planning a trip, or
    saving for a major expense.
    
    \item \textbf{Irregular income and expenses:} traditional monthly budgets
    often assume predictable income and expenses and therefore provide limited
    guidance when financial conditions vary.
    
    \item \textbf{Weak behavioral feedback:} users may repeatedly exhibit
    spending patterns without receiving specific, timely, and actionable
    suggestions.
    
    \item \textbf{Shared financial complexity:} households may have multiple
    contributors, shared expenses, different ownership arrangements, and
    approval requirements that are difficult to represent using simple
    expense-splitting mechanisms.
    
    \item \textbf{Limited planning under uncertainty:} a single projection may
    give users a false sense of certainty when future income, expenses, and
    savings behavior can change.
    
    \item \textbf{Lack of transparent intervention:} financial applications
    may recommend actions without clearly explaining why a recommendation was
    generated or allowing users to inspect and modify the underlying rules.
\end{itemize}

\textbf{Why this matters now (contextual relevance):} Users increasingly manage
multiple financial goals and irregular expenses while making decisions under
uncertainty. A system that connects financial behavior with future goals can
provide more useful guidance than a system that only reports historical
transactions.


\section{Proposed Solution%
  \BvrHint{\ldots Software Proposition}}

\subsection{Overview}

We propose a \textbf{web-based Smart Personal Finance \& Goal Management
System} that combines financial tracking, goal-based budgeting, behavioral
analysis, scenario simulation, and actionable interventions.

The major components are:

\begin{itemize}[leftmargin=0.7cm]
    \item \textbf{Financial tracking:} users record and categorize income,
    expenses, accounts, and transactions.
    
    \item \textbf{Goal management:} users define financial goals with target
    amounts, deadlines, priorities, and progress.
    
    \item \textbf{Goal-based budgeting:} the system connects spending and
    savings behavior with financial goals and provides recommendations for
    staying on track.
    
    \item \textbf{Scenario simulation:} users can compare different financial
    scenarios and use Monte-Carlo-style simulations to estimate the likelihood
    of achieving a goal under varying income, expenses, and savings conditions.
    
    \item \textbf{Habit and transaction analysis:} the system detects recurring
    spending patterns and links them to user-defined or automatically
    suggested financial rules.
    
    \item \textbf{Anomaly detection and intervention:} unusual or rapidly
    increasing spending patterns are identified and converted into gentle,
    concrete recommendations rather than only displaying warnings.
    
    \item \textbf{Shared household finances:} users can manage shared expenses,
    ownership, approval flows, contributions, and fairness metrics.
    
    \item \textbf{Audit and decision support:} users can export financial
    records and compare scenarios for significant financial decisions.
\end{itemize}


\subsection{Core workflow}

\begin{enumerate}[leftmargin=0.7cm]
    \item A user securely creates an account and records or imports financial
    transactions.
    
    \item Transactions are categorized and associated with relevant accounts,
    goals, budgets, or household members.
    
    \item The system analyzes historical financial behavior to identify
    recurring patterns and unusual changes.
    
    \item The user creates a financial goal by specifying a target amount,
    deadline, and other relevant preferences.
    
    \item The system estimates whether the goal is achievable using current
    financial behavior and generates alternative scenarios.
    
    \item The user can modify assumptions such as income, expenses, or monthly
    savings and compare the resulting scenarios.
    
    \item When the system detects a potentially important spending pattern or
    deviation, it provides an explainable intervention with a concrete
    suggested adjustment.
    
    \item Users can inspect, modify, accept, or reject the rules behind
    automated recommendations.
\end{enumerate}


\subsection{Operational constraints for a privacy-first financial system}

\begin{itemize}[leftmargin=0.7cm]
    \item Financial information must be protected through secure authentication,
    authorization, encryption, and least-privilege access.
    
    \item The system should minimize collection of sensitive information and
    store only the data required for its intended functionality.
    
    \item Recommendations should be explainable and should not make
    irreversible financial decisions automatically.
    
    \item The system will provide financial planning and visualization features
    without attempting to function as a full brokerage platform.
    
    \item Money calculations must be implemented carefully and tested
    extensively to avoid rounding and arithmetic errors.
\end{itemize}


\section{Solution Approach%
  \BvrHint{\ldots Engineering focus}}

This is an engineering course project, so the focus will be on reliable system
design, secure handling of financial information, explainable decision
support, testing, and maintainability rather than purely research-oriented
algorithmic novelty.


\subsection{Goal-based budgeting and scenario simulation}

The system will connect current financial behavior with future goals. Instead
of assuming a single deterministic future, the system will allow users to
compare multiple scenarios.

Monte-Carlo-style simulation will be used to generate multiple possible future
financial paths based on assumptions about income, expenses, savings, and other
relevant variables. The resulting distribution can be used to estimate the
likelihood of achieving a financial goal by a specified deadline.

Users will be able to compare scenarios such as:

\begin{itemize}[leftmargin=0.7cm]
    \item maintaining current spending behavior,
    \item reducing selected discretionary expenses,
    \item increasing monthly savings,
    \item experiencing lower or irregular income, and
    \item changing the target date or goal amount.
\end{itemize}


\subsection{Habit and transaction rule engine}

The system will analyze transaction history to identify recurring or
meaningful spending patterns. Based on these patterns, it can suggest rules
that users can inspect and modify.

For example, repeated spending in a particular category may result in a
suggestion to set a monthly limit or redirect a portion of future savings
toward a selected goal.

The rule engine will be designed to remain transparent so that users can
understand why a recommendation was generated and modify or disable it when
required.


\subsection{Anomaly detection and intervention}

The system will identify unusual spending behavior using historical
transaction patterns and category-level trends.

Instead of presenting only an ``anomaly detected'' message, the system will
provide a gentle intervention containing:

\begin{itemize}[leftmargin=0.7cm]
    \item the category or behavior that changed,
    \item an explanation of the deviation,
    \item the potential impact on relevant goals, and
    \item a concrete suggested adjustment.
\end{itemize}


\subsection{Shared household finance workflow}

The system will support multiple users participating in shared financial
activities. Each shared transaction can have ownership information and,
where appropriate, an approval state.

The system will maintain contribution information and provide fairness metrics
based on configurable sharing rules rather than assuming that all expenses
must always be divided equally.


\subsection{Web architecture%
  \BvrHint{\ldots web-based solution}}

A typical 3-tier architecture will be used:

\begin{itemize}[leftmargin=0.7cm]
    \item \textbf{Frontend:} responsive user interface for dashboards,
    transactions, budgets, goals, scenarios, and household workflows.
    
    \item \textbf{Backend API:} authentication, authorization, transaction
    processing, budgeting logic, rule execution, anomaly detection, scenario
    simulation, and audit logging.
    
    \item \textbf{Database:} storage for users, accounts, transactions,
    categories, goals, budgets, rules, scenarios, household relationships,
    approvals, and audit records.
\end{itemize}


\subsection{Security and auditability}

Financial information will be treated as sensitive application data. The
system will incorporate:

\begin{itemize}[leftmargin=0.7cm]
    \item authentication and role-based authorization,
    \item least-privilege access,
    \item encryption for sensitive data and secure communication,
    \item secure secret management,
    \item validation of financial inputs, and
    \item append-oriented audit logs for important financial actions.
\end{itemize}


\subsection{CI/CD and testing}

CI/CD will be used to automatically build the application, run automated
tests, perform static checks, and produce repeatable deployment artifacts.

Special attention will be given to financial calculations. Unit tests,
integration tests, and property-based tests will be used to validate important
properties of money calculations, budgeting logic, scenario calculations, and
goal projections.


\section{Evaluation Criterion%
  \BvrHint{\ldots measurable and attributable}}

Success will be evaluated using measurable criteria related to correctness,
goal planning, recommendation quality, and system reliability.


\subsection{Primary evaluation metric}

\textbf{Goal projection accuracy and consistency:} the system should produce
consistent and mathematically valid goal projections when identical assumptions
are supplied.

The evaluation will compare simulated outcomes against controlled test
scenarios with known expected behavior.


\subsection{Secondary evaluation metrics}

\begin{itemize}[leftmargin=0.7cm]
    \item \textbf{Goal achievement feasibility:} percentage of test scenarios
    in which the system correctly identifies whether a goal is feasible under
    the specified assumptions.
    
    \item \textbf{Recommendation usefulness:} user evaluation of whether
    suggested interventions are understandable, relevant, and actionable.
    
    \item \textbf{Anomaly detection quality:} ability to identify meaningful
    deviations without producing excessive false alerts.
    
    \item \textbf{Calculation correctness:} correctness of balances, budgets,
    contributions, projections, and other money-related calculations.
    
    \item \textbf{System reliability:} successful completion of critical
    workflows such as authentication, transaction recording, budgeting, and
    goal creation.
    
    \item \textbf{Security and auditability:} verification that unauthorized
    users cannot access protected financial information and that important
    financial actions generate appropriate audit records.
\end{itemize}


\subsection{Pilot validation plan%
  \BvrHint{\ldots high-level}}

A small group of users will be given representative financial scenarios and
asked to complete common tasks such as creating a goal, categorizing
transactions, comparing scenarios, and responding to spending interventions.

The evaluation will combine automated test results with user feedback on
clarity, usefulness, and ease of understanding.


\section{Scalability%
  \BvrHint{\ldots with foresight}}

The system should be scalable both technically and functionally.

\begin{enumerate}[leftmargin=0.7cm]

    \item \textbf{Scalable architecture:} the frontend, backend services, and
    database will be separated so that individual components can evolve
    independently.
    
    \item \textbf{Database scalability:} indexes, pagination, efficient
    queries, and appropriate data modelling will be used for transaction-heavy
    workloads.
    
    \item \textbf{Time-series scalability:} transaction history and financial
    events will be stored in a form suitable for historical analysis and
    forecasting.
    
    \item \textbf{Simulation scalability:} scenario simulations can be
    separated from regular request processing when computational requirements
    increase.
    
    \item \textbf{Operational scalability:} CI/CD and containerized or
    platform-hosted deployment can be used to support repeatable deployments.
    
\end{enumerate}


\section{Engine availability heuristic}

A mature set of software engineering tools and technologies is available to
implement the proposed system:

\begin{itemize}[leftmargin=0.7cm]
    \item Web framework for the frontend and backend application.
    
    \item Relational database for structured financial and user data.
    
    \item Authentication and authorization mechanisms for secure access.
    
    \item Time-series processing and forecasting libraries for financial
    history and scenario analysis.
    
    \item Statistical and numerical libraries for Monte-Carlo-style
    simulations.
    
    \item Automated testing frameworks including support for property-based
    testing.
    
    \item CI/CD tools for automated testing and deployment.
\end{itemize}

These mature technologies will allow the project to focus on financial
workflow design, correctness, privacy, explainability, and user behavior
rather than building infrastructure from scratch.


\section{Project scope and deliverables%
  \BvrHint{\ldots iteration-friendly}}

\subsection{Initial deliverable%
  \BvrHint{\ldots first iteration}}

\begin{itemize}[leftmargin=0.7cm]
    \item Secure user authentication and basic authorization.
    
    \item Income and expense transaction management.
    
    \item Transaction categorization and financial dashboard.
    
    \item Creation and tracking of financial goals.
    
    \item Basic goal-based budgeting.
    
    \item Initial database schema and secure backend API.
    
    \item Automated unit and integration tests for core financial calculations.
    
    \item CI pipeline for building and testing the application.
\end{itemize}


\subsection{Subsequent deliverables}

\begin{itemize}[leftmargin=0.7cm]
    \item Habit and transaction pattern detection.
    
    \item User-inspectable rule engine for financial recommendations.
    
    \item Anomaly detection and gentle intervention system.
    
    \item Monte-Carlo-style goal projections and scenario comparison.
    
    \item Shared household finances with ownership and approval workflows.
    
    \item Fairness metrics for shared financial contributions.
    
    \item Tax-aware and investment-aware financial views without brokerage
    functionality.
    
    \item Exportable audit trail and decision scenario reports.
    
    \item Property-based testing for important financial calculations.
\end{itemize}


\section{Risks and mitigations}

\begin{itemize}[leftmargin=0.7cm]

    \item \textbf{Financial data privacy:} minimize sensitive data collection,
    enforce least-privilege access, use encryption, and maintain secure audit
    logs.
    
    \item \textbf{Incorrect financial calculations:} use precise numerical
    representations, clearly defined calculation rules, extensive automated
    testing, and property-based tests.
    
    \item \textbf{Unhelpful recommendations:} keep recommendations
    explainable, allow users to inspect and modify rules, and avoid
    automatically making irreversible financial decisions.
    
    \item \textbf{False anomaly alerts:} use historical patterns and
    configurable thresholds while allowing users to dismiss or provide
    feedback on alerts.
    
    \item \textbf{Simulation uncertainty:} clearly communicate that scenario
    projections are estimates based on assumptions rather than guaranteed
    outcomes.
    
    \item \textbf{Scope complexity:} implement the core tracking, goal, and
    budgeting functionality first and introduce advanced features in
    subsequent iterations.
    
    \item \textbf{Security vulnerabilities:} use secure coding practices,
    dependency checks, automated testing, and restricted access to sensitive
    resources.
    
\end{itemize}


\section{Summary}

This project proposes a privacy-first Smart Personal Finance \& Goal
Management System designed to move beyond passive financial tracking toward
goal-oriented decision support and behavioral change. The system will combine
transaction management, goal-based budgeting, behavioral pattern detection,
anomaly-based interventions, scenario simulation, shared household finances,
and an inspectable rule engine.

The project emphasizes software engineering depth through secure financial
data handling, scalable architecture, time-series analysis, forecasting,
Monte-Carlo-style simulation, explainable recommendations, auditability, and
strong automated testing of money-related calculations.

By connecting everyday financial behavior with long-term goals and showing
users the potential impact of alternative decisions, the system aims to help
users understand not only \emph{where their money went}, but also
\emph{what they can do next} to improve their financial outcomes.

\end{document}
